\subsection{Dual-use implications and security research boundary}\label{sec:dualuse}
\paragraph{Training-history inference is not unrestricted reverse engineering.}
ChronoTrace infers an order within a supplied catalogue of training maps; it does not infer an unknown architecture, recover undisclosed training data, extract weights from a hosted model, or assign semantic meanings to individual parameters. Its white-box access already includes the base state and executable gradient recipes. An API-only observer does not have these inputs, and weights alone do not provide the missing recipes or regularity guarantees. The present evidence therefore establishes neither a black-box extraction attack nor a general method for decoding what model weights mean.

\paragraph{A conditional route to studying intermediate states.}
Under the inverse assumptions of Lemma~\ref{lem:inverse}, a compatible suffix and endpoint observation induce an enclosure of its earlier state. Resolving chronology can narrow which such enclosure is relevant. This is a consequence of the existing transport bound, not a new experimental result about recovering earlier LLM checkpoints. A unique order need not yield a tight state estimate: Eq.~\eqref{eq:expansion} can amplify export and numerical uncertainty. Nor does an earlier state identify a selectively removable capability. A training stage changes parameters through a trajectory-dependent map; it is not, in general, an architectural layer or an additive block that can be removed while leaving later learning unchanged. Counterfactual omission of a stage, exact historical rollback, and semantic model editing are different tasks.

\paragraph{Safety relevance without an unsupported attack claim.}
Prior work shows that subsequent fine-tuning can degrade aligned language-model behaviour, even with benign data~\cite{qi2023safety}. Mechanistic studies also connect refusal behaviour to particular internal representations in tested models~\cite{arditi2024refusal}; this does not identify a universally detachable safety layer, nor equate refusal with all aspects of safety. Tamper-resistant safeguarding is an existing research problem for open-weight models~\cite{tamirisa2024tamper}. These findings motivate examining whether additional chronology information could aid tampering. ChronoTrace has not located safety-specific circuits, removed an alignment stage from an LLM, recovered a pre-alignment LLM checkpoint, or measured increased harmful capability. A verified stage order also does not prove that a stage achieved its intended safety objective or that the released model is safe.

\paragraph{The threat is additional information under existing access.}
A possible misuse is helping a sufficiently informed, weight-access adversary narrow a pipeline hypothesis or focus subsequent model analysis. This is a threat hypothesis, not a demonstrated capability gain. In the full-access regime, the base and recipes already permit substantial counterfactual replay, so attributing that pre-existing power to chronology inference would exaggerate the contribution to risk. A valid test must hold access and resources fixed, charge chronology acquisition, and compare with the same analysis performed without inferred order. Security depends on the additional information or cost reduction, not merely the existence of an inverse computation.

\paragraph{Defensive investigation and release boundaries.}
A proposed follow-up would use authorized small models, benign stage-specific properties, and held-out intermediate checkpoints to compare hidden-order, inferred-order, and evaluator-known-order conditions. Relevant outcomes are recovery error and enclosure coverage, cost at matched total work, unintended changes to unrelated benign behaviour, and robustness to recipe mismatch, unknown optimizer state, and export precision. This protocol has not been executed as part of the present study; it must not be counted as evidence of either an attack or a defense. Identifying a relevant stage alone would not demonstrate control of a behaviour. Safety-specific evaluation would require separate authorization, containment, and specialist review before any release of operational artifacts.

Signed training records, controlled access to checkpoints and recipes, and independent behavioural evaluation should complement chronology auditing rather than be replaced by it. A failure to reconstruct chronology is not evidence of tamper resistance, and a matching chronology is not a cryptographic attestation. The reference artifacts contain small-model chronology experiments, not a validated LLM safeguard-removal procedure. The dual-use question is left explicit: does chronology information add a measurable manipulation advantage beyond what the same-access observer could already do?

